\فصل{جداسازی مبتنی بر اطلاعات متقابل}
\label{algs}
در این فصل، الگوریتمی مبتنی بر کمنیه‌سازی اطلاعات متقابل، برای جداسازی برای مخلوط‌های غیرخطی را معرفی کرده و عملکرد آن‌ را بررسی می‌کنیم.

\section{تعاریف و مقدمات}

در ابتدا به بیان تعاریف و مقدمات لازم برای معرفی الگوریتم جداسازی می‌پردازیم. منبع ما در این قسمت، 
\cite{mbzadehthesis}
است.
\begin{definition}
اطلاعات متقابل متغیر‌های تصادفی
$x_1, x_2, \dots, x_N$
را به صورت زیر تعریف می‌کنیم:
\begin{equation}
I(\x) = \mathrm{KL}\Big(P_\x(\x) || \prod_{i = 1}^{N} p_{x_i}(x_i)\Big) = \sum_{i = 1}^{N} H(x_i) - H(\x)
\end{equation}
\end{definition}

از خواص فاصله‌ی 
\lr{KL}
بین دو توزیع، می دانیم که اطلاعات متقابل صفر است، اگر و تنها اگر داشته باشیم:
\begin{equation}
p_\x(\x) = \prod_{i = 1}^{N} p_{x_i}(x_i)
\end{equation}

\begin{definition}
\textbf{تابع امتیاز}
تابع امتیاز یک متغیر تصادفی 
$x$
با 
\lr{pdf}
$p_x(x)$
را به صورت زیر تعریف می‌کنیم:
\begin{equation}
\psi_x(x) \triangleq -\frac{d}{dx} \ln p_x(x) = -\frac{p_x'(x)}{p_x(x)}
\end{equation}
\end{definition}

\newpage
\begin{definition}
	برای یک بردار تصادفی 
$\x$، 
سه کمیت زیر را تعریف می‌کنیم:
\begin{itemize}
\item
\textbf{\lr{Marginal Score Function (MSF)}}
\begin{equation}
\mathbf{\psi}(\x) = (\psi_1(x_1), \psi_2(x_2),\dots, \psi_n(x_n))^\top
\end{equation}
که در آن
\begin{equation}
\psi_i(x_i) = - \frac{p'_{x_i}(x_i)}{p_{x_i}(x_i)}
\end{equation}
\item 
\textbf{\lr{Joint Score Function (JSF)}}
\begin{equation}
\mathbf{\varphi}(\x) = \Big(\varphi_1(x_1), \varphi_2(x_2),\dots, \varphi_n(x_n)\Big)^\top
\end{equation}
که در آن
\begin{equation}
\varphi_i(\x) = -\frac{\frac{\partial}{\partial x_i}p_\x(\x)}{p_\x(\x)}
\end{equation}

\item
\textbf{\lr{Score Function Difference (SFD)} }	
\begin{equation}
\mathbf{\beta}_\x (\x) \triangleq \mathbf{\psi}_\x(\x) - \varphi_\x(\x)
\end{equation}

\end{itemize}
\end{definition}

این کمیت‌ها ویژگی‌های بسیار جالبی دارند که در ادامه به آن‌ها اشاره خواهد شد:
\begin{theorem}
درایه‌‌های بردار تصادفی 
$\x$
مستقل هستند، اگر و تنها اگر 
$\mathbf{\beta}_\x = 0$
باشد.
\end{theorem}
\begin{proof}
	ساده است.
\end{proof}	

\begin{theorem}
اگر 
$\x$
یک بردار تصادفی با چگالی 
$p_\x$
و 
\lr{JSF}
برابر 
$\mathbf{\varphi_\x}$
بوده و تابع
$f(\x)$
تابعی با مشتقات جزئی مرتبه‌ی اول پیوسته باشد، اگر برای هر 
$i$
داشته باشیم:

\begin{equation}
\lim_{x_i \to \pm\infty} \int_{x_1, \dots, x_{i-1}, x_{i+1}, \dots, x_N} f(\x)p_\x(\x) dx_1 \dots dx_{i-1}dx_{i+1}\dots dx_N = 0
\end{equation}
آن‌گاه داریم:
\begin{equation}
\E\{f(\x)\phi_i(\x)\} = \E\Big\{\frac{\partial f}{\partial x_i}(\x)\Big\}
\end{equation}
\end{theorem}
\begin{proof}
	به سادگی با انتگرال‌گیری جزء به جزء به دست می‌آید.
\end{proof}

\begin{theorem}
	فرض کنید 
$\x$
یک بردار تصادفی کران‌دار باشد و 
$\Delta$
نیز یک بردار تصادفی کوچک. در این صورت:
\begin{equation}
I(\x + \Delta) - I(\x) = \E\{\Delta^\top \mathbf{\beta}_\x(\x)\} + o(\Delta)
\end{equation}
\label{grad}
\end{theorem}

\begin{proof}
اثبات این قضیه، در پیوست B 
\cite{mbzadehthesis}
آمده است.
\end{proof}

\section{طراحی الگوریتم}
 قضیه
 \eqref{grad}‌، پیشنهاد طراحی الگوریتمی برای یافتن تابعی هموار است که با گذر دو فرآیند تصادفی داده‌شده از آن، اطلاعات متقابل آن کمینه ‌شود.
 
 شهوداً، قضیه‌ی 
 \eqref{grad}
 بیان می‌کند که 
 \lr{SFD}
 یک گرادیان «تصادفی» برای اطلاعات متقابل است و می‌توان از آن برای طراحی یک الگوریتم
 \lr{Steepest-Descent}
 بهره جست.  الگوریتم ما، تعمیمی از الگوریتم مطرح‌شده در
\cite{mbzadehthesis}
برای فرآیند‌های تصادفی است.  در 
\cite{mbzadehthesis}
 الگوریتم زیر ارائه شده است که در هر مرحله، منابع خروجی را در جهت «گرادیان» اطلاعات متقابل، به‌روز می‌کند.
\begin{figure}[h!]
\centering
\includegraphics[scale=0.5]{alg.png}
\end{figure}

این الگوریتم قابل تعمیم به فرآیند‌های تصادفی نیز هست. برای این کار، کافی است در هر مرحله، شیفتی تصادفی از یکی از فرآیند‌های تصادفی در نظر گرفت و الگوریتم فوق را بر روی این ورژن شیفت یافته انجام داد.
\newpage
\section{شبیه‌سازی عملکرد الگوریتم}
در این شبیه‌سازی، منابع دو فرآیند تصادفی
\lr{ARIMA}
مستقل
با 


\lr{
\texttt{model = arima('Constant',0.3,'AR',0.7,'MA',0.4,...\\
	'Distribution',tdist,'Variance',0.15)}}


ساخته و با تابع غیرخطی: 
\begin{equation}
\begin{cases}
x_1(t) = s_1(t) + 0.5s_2(t) + 0.5s_1(t)s_2(t)\\
x_2(t) = 0.6s_2(t) + 2s_2(t) + 0.3s_1(t)s_2(t)
\end{cases}
\end{equation}
مخلوط  می‌کنیم. الگوریتم را بر مخلوط‌ها اعمال کرده و مقدار
\lr{SNR}
جداسازی را بر حسب تعداد تکرار الگوریتم رسم می‌کنیم.
\begin{figure}[h!]
\centering
\includegraphics[scale=0.8]{snr.png}
\caption{
\lr{SNR}
جداسازی مخلوط غیرخطی با کمینه‌کردن اطلاعات متقابل
}
\end{figure}

این الگوریتم، قادر به رسیدن به 
\lr{SNR}
ای برابر با 
$15.5 dB$
و 
$18.0 dB$
است. در پروژه‌‌ی کارشناسی دو، عملکرد این الگوریتم با معدود الگوریتم‌های موجود برای جداسازی مخلوط‌های غیرخطی مقایسه خواهد شد. همچنین کیفیت عملکرد آن بر روی مسئله‌ی واقعی نیز بررسی خواهد گشت.
 